\documentclass[12pt, letterpaper]{article}

\title{Recursion \color{red}\textbf{[KEY]}\normalcolor}
\author{Amiel L. Iliesi}
\date{2026-03-27}

% packages
\usepackage[hidelinks]{hyperref}
\usepackage[x11names]{xcolor}
\usepackage{minted}

% new commands
\newcommand{\question}[1]{\textbf{#1}\par}
\newcommand{\answer}[1]{\color{Green4}\textbf{#1}\normalcolor}

% new environments
\newenvironment{answerblock}{\color{Green4}}{\normalcolor\vspace{1em}}

% global settings
\setlength{\parindent}{0pt}
\setlength{\parskip}{0.5em}
\usemintedstyle{monokai}

\begin{document}

\maketitle

\tableofcontents

\pagebreak
\section{Terminology}

\question{What are base cases?}
\begin{answerblock}
Base cases are similar to your iterative looping condition. They are kind of
inverse though, since your iterative condition tells you when to
\emph{keep looping}, whereas your base cases for recursion tell you when to
\emph{stop looping}.
\end{answerblock}

\question{What are recursive cases?}
\begin{answerblock}
These have no direct analog to iteration. They are your looping condition.
They are separate from base cases, because--in recursion--your reasons to go,
and reasons to stop, aren't necessarily the same.
\end{answerblock}

\question{What is recursive depth? Stack depth?}
\begin{answerblock}
Recursive depth is how many function calls deep you currently are. Stack-depth
is related and closely coupled; it is how many function calls on the underlying
stack you're in. Unless your compiler does something smart, they are usually
the same.
\end{answerblock}

\question{Describe the differences between preorder, in-order, and postorder
traversal. Also, which are synonyms with head recursion, and tail recursion?}

\begin{answerblock}
Preorder means you loop before \emph{(pre)} doing the work in your function.

In-order means that--if your data structure has an ordering--you recurse on
anything that comes before you, then when you are about to execute the current
node relative to the ordering \emph{(in-order)}, you do your work, and finally
continue to the rest of the connections. This doesn't apply to recursion with
lists.

Postorder means you loop after \emph{(post)} doing all the work in your
function.

Tail recursion and postorder traveral are synonyms. Head recursion and preorder
traversal are synonyms.
\end{answerblock}

\begin{minipage}{\textwidth}
\question{When possible, why is \emph{tail recursion} preferable?}
\begin{answerblock}
Because tail recursion can be directly and systematically converted to
iteration, which your compiler can do for you automatically.
\end{answerblock}
\end{minipage}

\pagebreak
\section{Conversion From Iteration to Recrusion}

\question{Convert the following into a recursive function.}

\begin{minted}[style=bw]{cpp}
int sum = 0;
for (int i = 1; i < N; ++i) {
    sum += i;
}
\end{minted}

\begin{minted}[bgcolor=darkgray]{cpp}
int sum(int i, int N) {
    if (i > N) {
        return 0;
    }
    else if (i == N) {
        return i;
    }
    else {
        return i + sum(i+1, N);
    }
}
\end{minted}

\question{What parts can you systematically translate? Identify the
relationship these pieces dictate between recursion and iteration.}

\begin{answerblock}
    You can translate \emph{all parts}. Iteration can \emph{always} be
    converted to recursion.

    \begin{itemize}
        \item state variables \(\rightarrow\) parameters
        \item accumulators \(\rightarrow\) return values
        \item loop conditions \(\rightarrow\) base cases
    \end{itemize}
\end{answerblock}

\pagebreak
\section{Uniquely Recursive}

\question{Describe the relationship between recursion and iteration.}

\begin{answerblock}
In terms of capability, iteration is a subset of recursion. Recursion can do
everything iteration can, and more.
\end{answerblock}

\question{To expand on the previous question, what can recursion do that
iteration cannot?--and vice versa.}

\begin{answerblock}
There is nothing iteration can do that recursion cannot. However, recursion
has things it can do that recursion cannot.

Recursion is capable of both backtracking, and branching. Both of which are
impossible in iteration without auxiliary data structures. 
\end{answerblock}

\question{Why is iteration preferable when possible?}

\begin{answerblock}
Because iteration incurs no operational and memory overhead. IE it is
generally faster.

Less measurably, iteration is usually more readable, and always more simple to
debug.
\end{answerblock}

\question{What situations can have a iterative \textit{and} recursive solution,
but have the recursive solution be preferable?}

\begin{answerblock}
There are some situations where a recursive solution is a much more natural
fit. In that case, it might be worth the performance cost to have something
more readable for your coworkers or peers.
\end{answerblock}

\end{document}